\begin{textsample}{POS Dim 1 – vem\_ed\_32 – Score 53.00 – vem\_ed\_32/t167.txt}  \label{ex:f1_pos_030}
Jornada dos PCIs 2025: Abertura No dia 9 de outubro ocorreu a Jornada dos Projetos Colaborativos Internacionais 2025.
O webinário online e gratuito \textbf{é} realizado anualmente desde 2021, pela Coordenadoria Geral do Ensino Superior de Graduação (CGESG) do CPS.
A edição deste ano abordou o tema “Internacionalização, pesquisa e \textbf{extensão}”.
O objetivo da Jornada dos Projetos Colaborativos Internacionais \textbf{é} debater \textbf{questões} relativas aos Intercâmbios Virtuais, com \textbf{foco} nos Projetos Colaborativos Internacionais (PCIs) \textbf{desenvolvidos} nas Fatecs e em projetos do tipo COIL (Collaborative Online International Learning – Aprendizagem Colaborativa International realizada on-line) ou BRaVE (Brazilian Virtual Exchange), desenvolvidos em instituições \textbf{que} \textbf{utilizam} abordagens similares, independentemente da nomenclatura adotada.
A programação \textbf{contou} com palestrantes da Universidade Federal de Santa Catarina, Universidade Federal Fluminense e Universidad Nacional de Río Cuarto (Argentina).
E apresentações de Projetos Colaborativos Internacionais (PCIs) realizados nas Fatecs de Pompeia, Sorocaba, Santo André e Bebedouro.
A \textbf{quinta} edição da Jornada dos PCIs registrou 142 inscrições.
Na pesquisa realizada após o evento com os participantes, 86 \% dos respondentes afirmaram estar “muito satisfeitos” com a organização.
As gravações estão disponíveis na playlist do canal da CGESG no YouTube: https://bit.ly/4rpWGPD O \textbf{número} 27 de VEm resume a Jornada dos Projetos Colaborativos Internacionais 2024 e traz um \textbf{histórico} das edições \textbf{anteriores} do evento: Jornada dos Projetos Colaborativos Internacionais 2024 - https://tinyurl.com/3a7hcpr3 Após a abertura da Jornada dos PCIs 2025, realizada por Otávio de Moraes, chefe de gabinete da presidência do Centro Paula Souza, houve uma breve apresentação sobre a equipe de Apoio à Internacionalização da Divisão de Extensão e Pesquisa no Ensino Superior da CGESG e sobre a estruturação dos PCIs, seu \textbf{histórico} e crescimento.
Patrícia Patrício, responsável pela comunicação da equipe, e Osvaldo Succi Junior, coordenador dos PCIs, comentaram os \textbf{números} e a evolução dos projetos, envolvendo pesquisa e \textbf{extensão}.
Entre 2013 e 2025, \textbf{foram} realizados 750 PCIs envolvendo 57 Fatecs e 94 Instituições de Ensino Superior internacionais, engajando cerca de 15 mil fatecanos e 15 mil estudantes das instituições parceiras.
Por ano, \textbf{são} realizados cerca de 120 PCIs nas Fatecs, o \textbf{que} coloca o Centro Paula Souza entre as cinco maiores instituições do mundo em \textbf{número} de projetos de Intercâmbio Virtual do tipo COIL.
Palestras Luciane Stallivieri (Universidade Federal de Santa Catarina) proferiu a pa- lestra magistral “Pontes entre Internacionalização, Extensão e Pesquisa”.
Defendeu a internacionalização como eixo transversal \textbf{que} articula ensino, pesquisa e \textbf{extensão}. “Não \textbf{são} polos isolados, e sim um sistema integrado onde \textbf{intercâmbios} virtuais atuam como conectores \textbf{práticos} entre ensino, comuni- de e investigação”.
A professora apresentou definições de Internacionalização em Casa, Internacionalização do Currículo e Internacionalização da Educação Superior, Intercâmbio Virtual e COIL, entre outros conceitos relacionados.
E deu sugestões para \textbf{que} os Intercâmbios Virtuais ajudem a atingir objetivos extensionistas, gerem pesquisas e promovam internacionalização.
Cíntia Rabello (Universidade Federal Fluminense) relatou o “International Vir- tual Exchange Project e suas \textbf{contribuições} para a aprendizagem de língua in- glesa e letramentos digitais de licenciandos na Universidade Federal Fluminense”.
Trata-se de uma iniciativa financiada pelo governo do Japão desde 2015.
A professora participou com seus alunos de inglês na UFF, em cinco edições do projeto, entre 2019 e 2025. “\textbf{É} uma oportunidade de \textbf{expandir} a aprendizagem de língua inglesa para além da sala de aula por meio da comunicação autêntica entre estudantes de diferentes culturas”, ressalta Cíntia. “Apesar do engajamento e interesse, alguns alunos têm dificuldade de realizar as \textbf{interações} fora da sala de aula”, \textbf{observa}.
María Soledad Fontana e Silvia Cristina Bertolo (Universidad Nacional de Río Cuarto, Argentina) \textbf{compartilharam} uma experiência COIL trilíngue, entre 16 alunos de Português e Francês, orientados por quatro docentes da universidade argentina e 18 estudantes de Relações Internacionais da Unesp, supervisionados por dois professores.
O projeto “Restaurante internacional: jantares temáticos e eventos culturais no Brasil, Argentina e França” consistia na elaboração de um flyer em português, espanhol e francês com o menu de um evento multicultural relacionando os três países.
As \textbf{produções} \textbf{foram} \textbf{publicadas} no Padlet.
Apresentações Vânia Regina Alves de Souza, da Fatec Pompeia, discorreu sobre os PCIs realizados na unidade com Symbiosis Centre for Management Studies (SCMS / Índia), Yerevan State University (Armênia), Eastern Oregon University (EOU / EUA) e Bukidnon State University (Filipinas).
O projeto com a Índia envolve também a Fatec Americana, aborda temas como satisfação e motivação no trabalho e já está na \textbf{terceira} edição.
O PCI com Armênia ocorre desde 2023, engajando também as Fatecs Araçatuba, Bragança Paulista, Carapicuíba, Guaratinguetá, Jales, Sebrae e Tatuí.
Com Eastern Oregon University, o projeto na área de estudos comparativos sobre agricultura no Brasil e EUA ocorreu no primeiro semestre de 2025, com participação das Fatecs Pompeia e Mogi das Cruzes.
Os alunos dos EUA estudaram os biomas brasileiros e os alunos brasileiros pesquisaram as regiões norte-americanas.
O PCI com a universidade das Filipinas \textbf{começou} em outubro de 2025.
Como principais benefícios dessas experiências, Vânia citou a melhoria na \textbf{interação} com os alunos, a participação em eventos internacionais, escrita de artigos acadêmicos (por professores e alunos) e \textbf{desenvolvimento} de \textbf{competências} \textbf{socioemocionais} (soft skills).
Wanderley do Prado relatou sua primeira experiência com PCIs, \textbf{que} ocorreu entre Fatec São Roque e DUOC UC (Chile) no segundo semestre de 2024, sobre elaboração de anúncios \textbf{utilizando} memes.
No primeiro semestre de 2025, o professor \textbf{iniciou} um projeto, realizado em inglês, entre Fatec Sorocaba e UNAM, agora na segunda edição. \textbf{Cada} grupo \textbf{propõe} um \textbf{produto} ou serviço a \textbf{ser} oferecido no Brasil e no México, abordando Objetivos de Desenvolvimento \textbf{Sustentável} (ODS).
Entre as principais motivações dos alunos, Wanderley \textbf{identifica} a prática do \textbf{idioma} estrangeiro em um \textbf{ambiente} interativo e divertido, além de conhecer novas \textbf{pessoas} e culturas.
Para os professores, perceber \textbf{que} o interesse dos alunos pelas aulas aumentou, \textbf{publicar} artigos com os parceiros internacionais sobre a experiência dos PCIs e participar de projetos de \textbf{extensão} \textbf{institucional} por meio dos PCIs. \textbf{Quanto} ao \textbf{idioma}, o professor \textbf{observou} a \textbf{diferença} entre os alunos da Fatec São Roque, \textbf{que} têm aulas de inglês no currículo, e os de Sorocaba, \textbf{que} não têm (e acabaram buscando cursos fora da faculdade). “Quando o professor Osvaldo Succi me procurou perguntando sobre a possibi- lidade de um Projeto Colaborativo Internacional com a DUOC do Chile, eu tive a ideia de trocar informações: os chilenos nos apresentam os métodos de homologação veicular e eu trabalho com eles a sistemática de medições de potência, torque e emissões de poluentes.
O professor Orlando Salvo Junior, da Fatec Santo André, trabalha com essa área de inspeção veicular e homologação, \textbf{que} \textbf{é} um tema novo para nós.
Ler norma \textbf{é} uma \textbf{coisa}, ver acontecendo na prática \textbf{é} completamente diferente”.
Outra vantagem apresentada pelo coordenador \textbf{foi} incluir o PCI como projeto de \textbf{extensão} \textbf{institucional}. “Isso \textbf{foi} realizador”.
Além disso, Fróes \textbf{destacou} o engajamento dos alunos, \textbf{que} se ajudaram na pro- dução dos vídeos dos ensaios técnicos com os veículos.
A atividade \textbf{reduziu} reprovação por falta e nota na \textbf{disciplina}.
A nota média final subiu de 6,6 para 7,4 e a assiduidade melhorou 17 \%.
Selma de Fátima Grossi, coordenadora do curso de Logística, e Lígia de Gran- di, professora de espanhol na Fatec Bebedouro, falaram sobre a sinergia entre docentes do \textbf{idioma} espanhol e da logística em um PCI realizado com a Unimi- nuto (Colômbia).
O projeto está em sua quarta edição e aborda o tema “Logística empresarial e \textbf{sustentável}”, estudando aspectos relacionados à \textbf{produção} de café, banana, flores, carne e cana-de-açúcar, importantes na economia dos dois países.
A pri- meira edição (\textbf{início} de 2024) enfocou a temática da logística verde na cadeia de suprimentos.
No segundo semestre de 2024, a \textbf{questão} \textbf{foi} “como \textbf{reduzir} o \textbf{impacto} \textbf{ambiental}” desses \textbf{produtos} representativos.
No \textbf{início} de 2025, o tema \textbf{foi} “\textbf{impacto} \textbf{ambiental} do transporte terrestre” (rodoviário, \textbf{que} prevalece nos dois países).
No segundo semestre de 2025, a \textbf{questão} \textbf{é} “Do lixo ao consumo” (economia circular).
As professoras ressaltaram \textbf{que} o PCI \textbf{permitiu} conhecer aspectos do país vizinho, melhorar no \textbf{idioma} e desenvolver \textbf{empatia}.

% matched lemmas: ambiental, ambiente, anterior, cada, coisa, começar, compartilhar, competência, contar, contribuição, desenvolvimento, desenvolvir, destacar, diferença, disciplina, empatia, expandir, extensão, foco, histórico, identificar, idioma, impacto, iniciar, institucional, interação, intercâmbio, início, número, observar, permitir, pessoa, produto, produção, propor, prático, publicar, quanto, que, questão, quinto, reduzir, ser, socioemocional, sustentável, terceiro, utilizar
\end{textsample}
